\input{../doc-class-cours.tex}

\begin{document}


\section*{Comment mesurer la Terre ? }

Il y a plus de 2 000 ans, un savant grec nommé Ératosthène a imaginé une méthode  pour déterminer la taille de la Terre en observant le Soleil, la géométrie et la proportionnalité.

\subsection*{1 - L'histoire}

Dans la ville de Syène, aujourd'hui Assouan, au moment du solstice d'été, vers midi, le Soleil éclaire le fond d'un puits. Cela indique que les rayons du Soleil sont verticaux à cet endroit. Il n'y a pas d'ombre projetée sur les objets.  

À Alexandrie, au même moment, un bâton projette une ombre.

Ératosthène comprend que cette différence peut lui permettre de mesurer la Terre !
Son raisonnement repose sur trois idées :

\begin{itemize}[label={$\bullet$}]
  \item La Terre est une sphère.
  \item Les rayons du Soleil sont parallèles lorsqu'ils atteignent la Terre, car le Soleil est extrêmement éloigné.
  \item La différence entre les angles mesurés en deux lieux permet d'estimer la fraction du tour terrestre qui les sépare, à condition de connaître leur distance et leur position géographique.
\end{itemize} 


\begin{figure}[H]
  \centering
  \includegraphics[width=0.7\linewidth]{4-culture/eratosthene_mesure_terre.png}
\end{figure}

\subsection*{2 - L'exercice}

\begin{itemize}[label={$\bullet$}]
  \item Les deux villes sont supposées situées approximativement sur le même méridien.
  \item Distance entre deux villes : $800$ km.
  \item Un cercle complet mesure $360\degree$.
  \item Angle mesuré entre les verticales locales : $7,2\degree$.
\end{itemize} 

\hrule

\begin{enumerate}
  \item[1.] Combien de fois l'angle de $7,2\degree$ est-il contenu dans un angle de $360\degree$ ?
  \item[2.] Quelle fraction du cercle représente l'angle de $7,2\degree$ ?
  \item[3.] Si $800$ km correspondent à cette fraction du tour terrestre, calculer la circonférence de la Terre.
  \item[4.] Calculer le rayon de la Terre.
  \item[5.] Expliquez pourquoi les rayons du Soleil doivent être considérés comme parallèles pour simplifier le raisonnement.
\end{enumerate}

\end{document}
